\section{Preliminaries and Problem Setting}
\label{sec:preliminaries}

We fix notation for diffusion-model sampling, define the verifier and NFE budget, and state the noise trajectory search problem.

\paragraph{Diffusion models.}
A diffusion model generates a sample in $\mathbb{R}^d$ over $T$ discrete denoising steps, indexed $t = T, T-1, \dots, 1$. Write $x_t \in \mathbb{R}^d$ for the latent at step $t$ ($x_T$ is the initial Gaussian noise, $x_0$ the final sample). The Denoising Diffusion Probabilistic Modeling (DDPM) formulation~\citep{ho2020denoising} couples a forward process that perturbs data with Gaussian noise and a backward process that converts noise back to data. Setting $\alpha_t = 1-\beta_t$ and $\bar{\alpha}_t = \prod_{s=1}^{t}\alpha_s$ for a fixed noise schedule $\beta_t\in(0,1)$, samples emerge by iterating
\begin{equation}
  x_{t-1}
  = \tfrac{1}{\sqrt{1 - \beta_t}}\!\left(x_t - \tfrac{\beta_t}{\sqrt{1 - \bar{\alpha}_t}}\,\varepsilon_\theta(x_t, t)\right) + \tilde\sigma_t\,\varepsilon_t,
  \qquad \varepsilon_t\sim\mathcal{N}(0, I_d),
  \label{eq:ddpm-reverse}
\end{equation}
from $x_T\sim\mathcal{N}(0, I_d)$ to $x_0$, where $\varepsilon_\theta$ is a learned noise predictor, $\tilde\sigma_t$ is the (fixed) reverse-process noise scale, and $\varepsilon_t \in \mathbb{R}^d$ is the per-step injected Gaussian noise (with $I_d$ the $d{\times}d$ identity). Score-based~\citep{song2019generative} and flow-based samplers~\citep{lipman2023flow,liu2023flow} share the same structure: a $T$-step backward trajectory driven by per-step stochastic noise $\varepsilon_t$.

\paragraph{Inference-time noise search.}

\begin{definition}[Verifier]
\label{def:verifier}
Let $\mathcal{X} \subseteq \mathbb{R}^d$ denote the sample space (e.g., images or latents) and $\mathcal{C}$ the space of conditioning signals (e.g., text prompts or class labels). A \emph{verifier} $v : \mathcal{X} \times \mathcal{C} \to \mathbb{R}$ scores a generated sample $x \in \mathcal{X}$ against its conditioning signal $c \in \mathcal{C}$, returning a scalar quality measure. The verifier is fixed during search and treated as a black box; results hold for any bounded, query-efficient $v$.
\end{definition}

\begin{definition}[NFE budget]
\label{def:nfe}
A single \emph{function evaluation} (NFE) is one forward pass through the denoising network, the dominant cost in diffusion sampling. Standard sampling consumes exactly $T$ NFE; noise trajectory search spends the remaining $B - T$ NFE on local refinement, where $B \ge T$ is the total budget.
\end{definition}

\begin{definition}[Noise trajectory search problem]
\label{def:nts}
Given a pretrained diffusion model with one-step transition $F_\theta : \mathbb{R}^d \times \{1, \dots, T\} \times \mathbb{R}^d \to \mathbb{R}^d$ that maps $(x_t, t, \varepsilon_t)$ to $x_{t-1}$ (the abstraction of \cref{eq:ddpm-reverse}, formalized in \cref{eq:transition}), a conditioning signal $c \in \mathcal{C}$, a verifier $v(\cdot, c)$ (\cref{def:verifier}), and a total budget of $B$ function evaluations (\cref{def:nfe}), find a sequence $\{\varepsilon_t^*\}_{t=1}^T$ that maximizes the verifier score of the final sample $x_0$ subject to total NFE $\le B$.
\end{definition}

The difficulty is sequential: $\varepsilon_t$ shapes $x_{t-1}$, which determines which $\varepsilon_{t'}$ are effective at later $t' < t$. This motivates the two-level decomposition of \cref{sec:framework}.
